\section{Билет 13. Метод Ньютона решения задач с ограничениями}

\ticketq{Метод Ньютона решения задач с ограничениями.}

\paragraph{Постановка.}
\[
  \varphi(x) \to \min , \qquad x \in S ,
\]
$\varphi$ выпукла и дважды дифференцируема, $S$ — выпуклый компакт,
$H(x) = \nabla^2\varphi(x)$. Итерации $x_{k+1} = x_k + \alpha_k s_k$.

\textbf{Идея.} Это метод Ньютона, перенесённый на задачу с ограничениями: там,
где метод условного градиента (билет~12) минимизирует на $S$ \emph{линейную}
модель, здесь минимизируется \emph{квадратичная}.

\paragraph{Квадратичная модель и направление.}
\[
  \psi_k(x) = \nabla\T\varphi(x_k)(x - x_k)
            + \tfrac12 (x - x_k)\T H(x_k)(x - x_k) ,
\]
\[
  y_k = \argmin_{x \in S}\psi_k(x) , \qquad
  \eta_k = \psi_k(y_k) \le \psi_k(x_k) = 0 , \qquad
  s_k = y_k - x_k .
\]

\paragraph{Лемма (вспомогательные факты).}
\textit{1) $x_k + \alpha s_k \in S$ при $\alpha \in [0,1]$;
2) $\nabla\T\varphi(x_k)s_k \le \eta_k \le 0$;
3) при сильной выпуклости с константой $m$: $\eta_k \le -\tfrac{m}{2}\norm{s_k}^2$.}

\textsc{Док-во.}
1) $x_k + \alpha s_k = (1-\alpha)x_k + \alpha y_k$ — выпуклая комбинация точек
$S$.

2) $\eta_k = \nabla\T\varphi(x_k)s_k + \tfrac12 s_k\T H(x_k)s_k$, второе
слагаемое $\ge 0$ ($H \succeq 0$ по выпуклости) $\Rightarrow$
$\nabla\T\varphi(x_k)s_k \le \eta_k$; а $\eta_k \le \psi_k(x_k) = 0$.

3) $\psi_k$ сильно выпукла с той же $m$, $y_k$ — её минимум на $S$
$\Rightarrow$ $\nabla\T\psi_k(y_k)(x_k - y_k) \ge 0$, и
\[
  0 = \psi_k(x_k) \ \ge\ \psi_k(y_k)
    + \underbrace{\nabla\T\psi_k(y_k)(x_k-y_k)}_{\ge\,0}
    + \tfrac{m}{2}\norm{x_k - y_k}^2
  \ \ge\ \eta_k + \tfrac{m}{2}\norm{s_k}^2 . \ \blacksquare
\]

\paragraph{Выбор шага.}
$\alpha_0 = 1$, $\eps, \lambda \in (0,1)$; $\alpha_k = \alpha_0\lambda^{i_k}$ —
первое значение, удовлетворяющее условию дробления
\[
  \varphi(x_k + \alpha s_k) - \varphi(x_k) \ \le\ \alpha\,\eps\,\eta_k
\]
(либо наискорейший спуск вдоль $s_k$: $\min_{\alpha>0}\varphi(x_k+\alpha s_k)$).

\textbf{Лемма (дробление конечно).} $H$ непрерывен на компакте $\Rightarrow$
$\norm{H(x)} \le M$, откуда
\[
  \varphi(x_k + \alpha s_k) \le \varphi(x_k) + \alpha\eta_k
    + \tfrac{M}{2}\alpha^2\norm{s_k}^2 ,
\]
и условие выполнено, как только
\[
  \alpha\eta_k + \tfrac{M}{2}\alpha^2\norm{s_k}^2 \le \alpha\eps\eta_k
  \iff
  \alpha \le -\frac{2(1-\eps)\eta_k}{M\norm{s_k}^2} > 0 \quad (\eta_k < 0) .
\]
Значит дробление конечно, и принятый шаг отделён от нуля:
$\alpha_k \ge \min\{1,\ \tfrac{2\lambda(1-\eps)}{MC^2}\abs{\eta_k}\}$,
$C = \operatorname{diam}S$. $\blacksquare$

\paragraph{Критерий останова.}
$\eta_k = 0$ $\Rightarrow$ $x_k$ минимизирует $\psi_k$ на $S$, и условие
оптимальности $\nabla\T\psi_k(x_k)(x - x_k) = \nabla\T\varphi(x_k)(x-x_k) \ge 0$
$\forall x \in S$ для выпуклой $\varphi$ означает, что $x_k$ — решение.

\paragraph{Теоремы о сходимости.}
\textbf{Т1.} \textit{$\varphi(x_k)$ монотонно убывает и
$\lim_k \varphi(x_k) = \varphi(x_*) = \min_{x\in S}\varphi(x)$ при любом $x_0$.}

\textsc{Схема.} Из условия дробления
$\varphi(x_k) - \varphi(x_{k+1}) \ge \eps\alpha_k\abs{\eta_k} \ge 0$:
последовательность убывает, ограничена снизу ($S$ компакт) $\Rightarrow$
сходится $\Rightarrow$ $\alpha_k\abs{\eta_k} \to 0$, а с оценкой шага снизу
$\eta_k \to 0$. Далее для $x \in S$, $t \in [0,1]$ из $\psi_k(x_k+t(x-x_k)) \ge \eta_k$
следует $t\nabla\T\varphi(x_k)(x-x_k) \ge \eta_k - \tfrac{t^2}{2}MC^2$; при
$t = \sqrt{\abs{\eta_k}}$ получаем
$\nabla\T\varphi(x_k)(x_* - x_k) \ge -\sqrt{\abs{\eta_k}}\big(1+\tfrac{MC^2}{2}\big)$,
и по выпуклости
$\varphi(x_k) - \varphi(x_*) \le \sqrt{\abs{\eta_k}}\big(1+\tfrac{MC^2}{2}\big) \to 0$.
$\blacksquare$

\textbf{Т2.} \textit{Если $m\norm{y}^2 \le y\T H(x)y \le M\norm{y}^2$,
$0 < m \le M$, то $x_k \to x_*$ \textbf{сверхлинейно}.}

\textbf{Т3.} \textit{Если дополнительно $\norm{H(x)-H(y)} \le R\norm{x-y}$, то
скорость \textbf{квадратичная}.}

\textsc{Схема Т2--Т3.} Вблизи $x_*$ при $\eps < \tfrac12$ принимается полный шаг
$\alpha_k = 1$, то есть $x_{k+1} = y_k$. Складывая условия оптимальности
$\big(\nabla\varphi(x_k) + H_k(y_k - x_k)\big)\T(x_* - y_k) \ge 0$ и
$\nabla\T\varphi(x_*)(y_k - x_*) \ge 0$ и обозначая
$r_k = \nabla\varphi(x_k) - \nabla\varphi(x_*) - H_k(x_k - x_*)$:
\[
  m\norm{y_k - x_*}^2 \le \norm{r_k}\,\norm{y_k - x_*}
  \ \Rightarrow\ \norm{x_{k+1} - x_*} \le \frac{\norm{r_k}}{m} .
\]
Без липшицевости $H$: $\norm{r_k} = o(\norm{x_k - x_*})$ — сверхлинейная
сходимость; при $\norm{H(x)-H(y)} \le R\norm{x-y}$:
$\norm{r_k} \le \tfrac{R}{2}\norm{x_k-x_*}^2$ — квадратичная. $\blacksquare$

\paragraph{Алгоритм.}
\begin{csteps}
  \item Задать $x_0 \in S$, $\eps, \lambda \in (0,1)$, $k = 0$.
  \item Построить $\psi_k(x)$ по $\nabla\varphi(x_k)$ и $H(x_k)$.
  \item Решить $y_k = \argmin_{x\in S}\psi_k(x)$; найти
        $\eta_k = \psi_k(y_k)$, $s_k = y_k - x_k$.
  \item $\eta_k = 0$ — конец, $x_k$ оптимальна.
  \item Найти $\alpha_k$ дроблением от $\alpha_0 = 1$ (либо наискорейшим
        спуском вдоль $s_k$).
  \item $x_{k+1} = x_k + \alpha_k s_k$, $k := k+1$, к~шагу~2.
\end{csteps}

\textbf{Замечание.} Для задачи с ограничениями-\emph{равенствами} $h(x) = 0$
тот же метод Ньютона записывают иначе: условия Лагранжа
$\nabla f + J\T\lambda = 0$, $h = 0$ рассматривают как систему $F(x,\lambda) = 0$
и применяют к ней метод Ньютона, получая на каждом шаге блочную систему ККТ
$\big(\begin{smallmatrix} W_k & J_k\T \\ J_k & 0\end{smallmatrix}\big)
\big(\begin{smallmatrix} p_k \\ \Delta\lambda_k\end{smallmatrix}\big)
= -\big(\begin{smallmatrix} \nabla_x L \\ h \end{smallmatrix}\big)$.

\seefull{гл.~3, \S10 (а вариант с системой ККТ — \S13)}
